	\section{Sección 2: Desarrollo del Fundamento Teórico de Ingeniería de Requisitos (Paso 3)}
\subsection{2.1 Fundamentos de la Ingeniería de Requisitos (Desarrollado por Jostin G. Vélez)}
\subsubsection{Concepto 1: Ingeniería de Requisitos como Disciplina Mediadora}
\textbf{Definición Formal:} La ingeniería de requisitos (IR) es una función interdisciplinaria que media entre los dominios del adquirente y el proveedor para establecer y mantener los requisitos que debe cumplir el sistema, software o servicio de interés \cite{ieee29148}. Se ocupa de descubrir, elicitar, desarrollar, analizar, verificar, validar, comunicar, documentar y gestionar dichos requisitos a lo largo de todo el ciclo de vida del sistema \cite{ieee29148, ieee15288}. Asimismo, un requisito se define formalmente como una propiedad o capacidad que debe exhibirse para resolver un problema del mundo real o alcanzar un objetivo específico de un stakeholder \cite{pohl2025, swebok2024}.
\textbf{Explicación Propia:} La ingeniería de requisitos no consiste simplemente en actuar como un transcriptor pasivo de lo que el cliente solicita inicialmente, sino que constituye un proceso analítico, sistemático y de ingeniería profundo. Implica investigar la raíz del problema organizacional subyacente y transformar intenciones humanas vagas, incompletas y frecuentemente contradictorias en especificaciones técnicas precisas, medibles y verificables sin atarlas prematuramente a decisiones de implementación. Cuando esta disciplina se ejecuta de manera superficial, los defectos analíticos se propagan por efecto cascada hacia el diseño y la construcción, multiplicando exponencialmente los costos de corrección y el riesgo de fracaso del proyecto.
\textbf{Ejemplo Vinculado al PFC:} En la Plataforma de Ventas para locales eléctricos de Quevedo, el síntoma inicial expresado por el dueño era ``la facturación es demasiado lenta y perdemos clientes''. Aplicar la ingeniería de requisitos permitió trascender la simple modificación de pantallas y descubrir la causa raíz: los cajeros demoraban hasta 5 minutos consultando catálogos físicos por desconocimiento electrotécnico. Esta necesidad se transformó en el requisito formal de contar con un módulo de recomendación inteligente que entregue tres sugerencias de productos técnicamente compatibles en un tiempo menor o igual a 3.0 segundos.
\subsubsection{Concepto 2: Articulación con el Ciclo de Vida de Sistemas y Software (ISO 15288 / ISO 12207)}
\textbf{Definición Formal:} Según los estándares ISO/IEC/IEEE 15288:2023 \cite{ieee15288} e ISO/IEC/IEEE 12207:2017 \cite{ieee12207}, los procesos técnicos de definición de necesidades de stakeholders y análisis de requisitos operan de forma recursiva e iterativa a través de los niveles de negocio (BRS), stakeholder (StRS), sistema (SyRS) y software (SRS) \cite{ieee29148}. Cada nivel de abstracción transforma necesidades operativas del entorno en requisitos verificables que alimentan la posterior definición arquitectónica \cite{ieee12207, hull2011}.
\textbf{Explicación Propia:} Los requisitos no nacen aislados; responden a una jerarquía de valor organizacional estricta. Una regla estratégica de negocio origina expectativas en los usuarios operativos, las cuales a su vez imponen capacidades al sistema sociotécnico global y, finalmente, se detallan en especificaciones funcionales y de calidad del software. Mantener esta alineación vertical garantiza que cada capacidad especificada aporte valor directo al objetivo empresarial y que ninguna función sea definida sin justificación operativa o económica.
\textbf{Ejemplo Vinculado al PFC:} En el PFC, el objetivo de negocio de ``incrementar las ventas cruzadas de material eléctrico en un 20\% trimestral'' (nivel BRS) se descompone en la necesidad del cajero de ``visualizar accesorios complementarios al facturar un componente principal'' (nivel StRS), lo que deriva en el requisito de software (nivel SRS) donde la Plataforma de Ventas identifica automáticamente rieles de montaje y terminales compatibles cada vez que el cajero registra un interruptor termomagnético en el comprobante.
\subsection{2.2 Categorías y Niveles de Requisitos (Desarrollado por Jostin G. Vélez)}
\subsubsection{Concepto 1: Requisitos Funcionales vs. Requisitos No Funcionales}
\textbf{Definición Formal:} De acuerdo con el SWEBOK v4.0 \cite{swebok2024} y Glinz \cite{glinz2007nfr}, los requisitos de software se categorizan fundamentalmente en funcionales y no funcionales. Los requisitos funcionales (RF) especifican los comportamientos observables, transformaciones de datos, políticas comerciales y reacciones que el software debe proporcionar ante estímulos de entrada \cite{pohl2025, swebok2024}. Por su parte, los requisitos no funcionales (RNF) actúan como atributos de calidad o restricciones que especifican niveles cuantificables de desempeño, seguridad, fiabilidad y usabilidad \cite{iso25010, glinz2007nfr}.
\textbf{Explicación Propia:} Separar rigurosamente los requisitos funcionales de los no funcionales (atributos de calidad) es indispensable para estructurar correctamente la especificación del sistema. Mientras los requisitos funcionales responden al ``qué'' hará el sistema para interactuar con los actores de negocio, los requisitos no funcionales describen ``qué tan bien'' lo hará, imponiendo límites medibles de tiempo de respuesta, seguridad y disponibilidad. Utilizar el filtro conceptual de la ``tecnología perfecta'' ayuda a discernir si un requisito existiría incluso con recursos computacionales infinitos, identificándolo así como funcional puro.
\textbf{Ejemplo Vinculado al PFC:} Un requisito funcional claro en nuestro PFC establece que la Plataforma de Ventas calcule el subtotal, el desglose del impuesto al valor agregado vigente (15\%) y el total a pagar cuando el cajero registre ítems eléctricos en el comprobante. Por el contrario, el requisito no funcional (calidad de servicio) exige que el sistema procese y confirme el registro completo de la venta y el descuento de inventario en menos de 250 milisegundos para el 95\% de las transacciones bajo una carga concurrente de 10 terminales de caja activas.
\subsubsection{Concepto 2: Requisitos de Negocio, de Stakeholder y de Solución}
\textbf{Definición Formal:} La norma ISO/IEC/IEEE 29148:2018 \cite{ieee29148} y Wiegers \& Beatty \cite{wiegers2013} estructuran los requisitos en tres estratos: (i) Requisitos de Negocio, que declaran el propósito de alto nivel de la organización; (ii) Requisitos de Stakeholder o Usuario, que describen las metas y tareas que los actores deben poder realizar con el sistema; y (iii) Requisitos de Solución (Sistema y Software), que detallan las capacidades funcionales, de calidad e interfaces necesarias para resolver el problema sin prescribir el diseño interno \cite{ieee29148, ireb2023}.
\textbf{Explicación Propia:} Confundir las metas de un usuario con las reglas financieras del negocio o con mecanismos prematuros de programación es una de las principales causas de ambigüedad en los proyectos de software. Estructurar los requisitos en estos tres estratos permite comunicarse con el propietario del local en términos de rentabilidad, con el cajero en términos de tareas operativas ágiles, y con los ingenieros de software mediante especificaciones precisas y verificables. Así se asegura una trazabilidad limpia desde la visión corporativa hasta las pruebas de aceptación.
\textbf{Ejemplo Vinculado al PFC:} En nuestra plataforma, el requisito de negocio exige ``eliminar las sanciones tributarias por emisiones extemporáneas de comprobantes''; el requisito de stakeholder del administrador establece ``poder auditar el estado de autorización de todos los comprobantes electrónicos emitidos en la jornada desde una vista consolidada''; y el requisito de solución de software especifica ``generar, firmar electrónicamente y transmitir el comprobante electrónico estructurado hacia el servicio de autorización tributaria del SRI inmediatamente al confirmar el cobro''.
\subsection{2.3 Taxonomía Detallada de Requisitos (Desarrollado por José I. Gilces)}
\subsubsection{Concepto 1: Requisitos de Interfaz, de Datos, de Rendimiento, de Conformidad y Restricciones}
\textbf{Definición Formal:} La norma ISO/IEC/IEEE 29148:2018 (cláusula 5.2.8) \cite{ieee29148} y el estándar de calidad ISO/IEC 25010:2023 \cite{iso25010} establecen una taxonomía granular para clasificar las especificaciones de un sistema:
\begin{enumerate}
	\item \textbf{Requisitos de Interfaz Externa:} Definen los límites compartidos entre el sistema y elementos externos, subdividiéndose en interfaces de usuario (UI), hardware (HW), software (SW) y comunicaciones (COM) \cite{ieee29148}.
	\item \textbf{Requisitos de Rendimiento (\textit{Performance Efficiency}):} Cuantifican el comportamiento temporal (latencia, tiempo de respuesta, capacidad de procesamiento) bajo cargas determinadas \cite{iso25010}.
	\item \textbf{Restricciones Operativas y de Dominio (\textit{Constraints}):} Limitaciones impuestas por el entorno organizacional, comercial o físico que acotan el espacio de solución \cite{pohl2025, swebok2024}.
	\item \textbf{Requisitos de Conformidad Regulatoria y Legal:} Mandatos derivados de leyes gubernamentales, normativas fiscales o estándares técnicos obligatorios \cite{ieee29148, ireb2023}.
\end{enumerate}
\textbf{Explicación Propia:} Una taxonomía detallada evita que el documento de especificación de requisitos (SRS) se convierta en un listado desordenado de peticiones mezcladas. Clasificar explícitamente los requisitos de interfaz permite definir qué flujos de información entran y salen por el límite del sistema hacia periféricos o entidades externas sin invadir el diseño interno del software. Del mismo modo, aislar las restricciones operativas y de conformidad legal garantiza que la solución respete el marco tributario ecuatoriano y las condiciones reales de trabajo del establecimiento comercial.
\textbf{Ejemplo Vinculado al PFC:} En el PFC se aplican todas estas subcategorías libres de implementación: como \textit{Interfaz de Hardware}, la captura de códigos de barras estándar EAN-13 desde lectores ópticos de mostrador; como \textit{Interfaz de Software y Comunicaciones}, el intercambio seguro de consultas técnicas y respuestas estructuradas con el servicio externo de inferencia de IA; como \textit{Restricción Operativa}, la ejecución íntegra desde navegadores estándar sin instalar software propietario en las cajas; y como \textit{Conformidad Legal}, el cumplimiento de la Ficha Técnica de Comprobantes Electrónicos del SRI y la Ley Orgánica de Protección de Datos Personales (LOPDP).
\subsubsection{Concepto 2: Sintaxis Estructurada EARS y Priorización MoSCoW}
\textbf{Definición Formal:} El enfoque EARS (\textit{Easy Approach to Requirements Syntax}), formulado por Mavin et al. \cite{mavin2009ears}, reduce la ambigüedad del lenguaje natural mediante cinco patrones sintácticos estructurados: \textit{Ubiquitous} (``The <system> shall <response>''), \textit{Event-driven} (``WHEN <trigger>, the <system> shall <response>''), \textit{State-driven} (``WHILE <state>, the <system> shall <response>''), \textit{Unwanted behavior} (``IF <condition>, THEN the <system> shall <response>'') y \textit{Optional} (``WHERE <feature>, the <system> shall <response>'') \cite{mavin2009ears}. Complementariamente, la técnica MoSCoW (\textit{Must have, Should have, Could have, Won't have}) clasifica la prioridad de entrega en función del valor crítico y el riesgo operativo \cite{ireb2023, wiegers2013}.
\textbf{Explicación Propia:} Redactar requisitos en prosa libre suele introducir verbos pasivos, cláusulas condicionales confusas y lagunas sobre qué evento dispara cada acción del sistema. Adoptar los patrones de EARS impone una disciplina gramatical rigurosa donde cada requisito funcional declara explícitamente su precondición o evento disparador, el sistema responsable y la respuesta observable esperada. Al combinar esta precisión sintáctica con la priorización MoSCoW, el equipo garantiza que el proyecto concentre sus esfuerzos en las capacidades vitales (\textit{Must}) sin descuidar la planificación evolutiva de mejoras deseables (\textit{Should/Could}).
\textbf{Ejemplo Vinculado al PFC:} En lugar de redactar vagamente ``el sistema debe manejar fallos de internet al usar la IA'', aplicamos el patrón de comportamiento no deseado (\textit{Unwanted behavior}) de EARS: ``\textit{IF the external AI recommendation service fails to respond within 3.0 seconds, THEN the Sales Platform shall display a local fallback recommendation list based on predefined electrical category rules within 200 ms}''. Asimismo, la emisión de comprobantes electrónicos SRI se prioriza como \textit{Must Have}, mientras que la personalización de temas visuales se clasifica como \textit{Could Have}.
\subsection{2.4 Propiedades de Calidad de los Requisitos según IEEE 29148 (Desarrollado por José I. Gilces)}
\subsubsection{Concepto 1: Propiedades de Calidad de un Requisito Individual (Énfasis en Apropiado al Nivel y Verificabilidad)}
\textbf{Definición Formal:} La norma ISO/IEC/IEEE 29148:2018 (cláusula 5.2.5) estipula que toda declaración de requisito individual bien formada debe satisfacer nueve características esenciales: (1) \textbf{Necesario} (su omisión provoca una deficiencia no mitigada), (2) \textbf{Apropiado al nivel / Libre de implementación} (describe el ``qué'' sin imponer innecesariamente el ``cómo'' ni tecnologías de diseño), (3) \textbf{Sin ambigüedad / Inambiguo} (admite una única interpretación semántica), (4) \textbf{Completo} (describe totalmente la capacidad necesaria), (5) \textbf{Singular} (expresa una sola función o atributo sin conjunciones múltiples), (6) \textbf{Factible} (realizable dentro de restricciones técnicas, de costo y plazo), (7) \textbf{Verificable} (susceptible de ser comprobado objetivamente mediante inspección, análisis, demostración o prueba cuantitativa), (8) \textbf{Correcto} (representa fielmente la necesidad real del stakeholder) y (9) \textbf{Conforme} (respeta el estilo y plantilla estándar, como EARS) \cite{ieee29148, pohl2025}.
\textbf{Explicación Propia:} Evaluar cada requisito individual contra estas nueve propiedades constituye el primer filtro de aseguramiento de calidad en la ingeniería de requerimientos. Una violación frecuente en la práctica académica consiste en vulnerar la propiedad de ser \textit{Apropiado al nivel} al incluir nombres de lenguajes de programación, motores de bases de datos o librerías específicas dentro del enunciado del requisito, lo cual contamina la especificación con diseño prematuro. Asimismo, erradicar adjetivos subjetivos como ``rápido'' o ``amigable'' mediante umbrales cuantitativos asegura que cada requisito sea objetivamente verificable en las pruebas de aceptación.
\textbf{Ejemplo Vinculado al PFC:} Si escribiéramos ``El sistema buscará productos rápidamente en la base de datos y exportará reportes amigables'', violaríamos las propiedades de \textit{Apropiado al nivel}, \textit{Singularidad} y \textit{Verificabilidad}. Corrigiéndolo bajo IEEE 29148 de forma agnóstica a la tecnología, lo expresamos como: ``\textit{WHEN the cashier enters at least 3 alphanumeric characters in the search field, the Sales Platform shall return matching electrical products from the catalog in less than 250 ms for inventories up to 50,000 items}''.
\subsubsection{Concepto 2: Propiedades de Calidad del Conjunto de Requisitos (SRS Completa)}
\textbf{Definición Formal:} Además de la calidad individual, ISO/IEC/IEEE 29148:2018 (cláusula 5.2.6) define las características que debe exhibir el conjunto integrado de requisitos dentro de una especificación (SRS): (1) \textbf{Completo} (contiene todos los requisitos relevantes, escenarios operativos e interfaces sin cláusulas pendientes ``TBD''), (2) \textbf{Consistente} (libre de contradicciones lógicas, temporales o terminológicas entre requisitos), (3) \textbf{Asequible / Factible en conjunto} (el sistema completo puede realizarse dentro del presupuesto y plazo global), (4) \textbf{Acotado} (mantiene el alcance alineado a los objetivos de negocio sin \textit{scope creep}) y (5) \textbf{Trazable} (tanto hacia atrás —origen/stakeholder— como hacia adelante —diseño posterior y pruebas—) \cite{ieee29148, swebok2024}.
\textbf{Explicación Propia:} Es perfectamente posible tener decenas de requisitos individualmente bien redactados que, al juntarse en un mismo documento, colisionen entre sí o desborden el presupuesto del proyecto. La consistencia del conjunto asegura que una exigencia de seguridad no bloquee un requerimiento de agilidad operativa en mostrador, mientras que la trazabilidad bidireccional permite realizar análisis de impacto precisos cuando un interesado solicita un cambio. Sin trazabilidad, cualquier modificación regulatoria o comercial obliga a revisar manualmente toda la especificación, elevando el riesgo de omisiones.
\textbf{Ejemplo Vinculado al PFC:} Durante el análisis del PFC detectamos un conflicto de \textit{Consistencia}: un requisito preliminar de seguridad exigía que el cajero ingresara su contraseña antes de cada consulta al módulo de recomendación IA, lo cual contradecía el requisito de rendimiento que exige obtener e insertar la recomendación en menos de 3.0 segundos durante la venta. Al auditar el conjunto bajo IEEE 29148, se resolvió exigir autenticación única por turno de caja (máximo 8 horas), preservando tanto la seguridad como la velocidad de atención y manteniendo la trazabilidad hacia el stakeholder Cajero (SH-01).
\subsection{2.5 Tipos de Sistemas y Contexto del Sistema (Desarrollado por Jeanpierre Robinson)}
\subsubsection{Concepto 1: Taxonomía de Sistemas Intensivos en Software e Híbridos con IA}
\textbf{Definición Formal:} Según Pohl (2025) \cite{pohl2025}, SWEBOK v4.0 \cite{swebok2024} e ISO/IEC/IEEE 15288:2023 \cite{ieee15288}, los sistemas intensivos en software se clasifican según su dominio y comportamiento en: (a) \textbf{Sistemas de Información Empresariales}, centrados en la persistencia, integridad transaccional y procesamiento de reglas de negocio; (b) \textbf{Sistemas Embebidos y Ciberfísicos}, estrechamente acoplados a sensores y actuadores de hardware con restricciones de tiempo real; y (c) \textbf{Sistemas Adaptativos o Basados en Inteligencia Artificial}, cuyo comportamiento incorpora modelos de inferencia probabilística que operan sobre contextos dinámicos \cite{pohl2025, swebok2024}.
\textbf{Explicación Propia:} Reconocer el tipo de sistema al que pertenece un proyecto determina qué criterios de calidad y restricciones deben priorizarse en la ingeniería de requisitos. Mientras un sistema de información tradicional requiere exactitud contable determinista, un sistema que incorpora capacidades de Inteligencia Artificial introduce latencias de procesamiento externo y el riesgo de generar respuestas imprecisas (\textit{alucinaciones}). Por consiguiente, cuando un sistema es híbrido (transaccional + IA), la especificación de requisitos debe establecer reglas de negocio estrictas para que el módulo inteligente jamás altere los cálculos financieros ni sugiera ítems inexistentes en el inventario real.
\textbf{Ejemplo Vinculado al PFC:} Nuestro PFC combina un núcleo transaccional determinista (facturación y control de inventario donde los cálculos monetarios e impositivos deben ser exactos al centavo) con un módulo adaptativo de recomendación basado en IA. Para gobernar esta dualidad desde los requisitos, especificamos que toda sugerencia generada por el módulo de IA debe estar restringida exclusivamente a productos registrados en el catálogo local que cuenten con saldo disponible mayor a cero en el almacén eléctrico de Quevedo.
\subsubsection{Concepto 2: Límite del Sistema, Zona Gris y las Ocho Perspectivas de Contexto de Pohl}
\textbf{Definición Formal:} Pohl (2025) \cite{pohl2025} e IREB CPRE v3.1 \cite{ireb2023} definen el \textbf{contexto del sistema} como la parte del entorno real que influye directamente en el sistema a desarrollar o que es afectada por él. El \textbf{límite del sistema} (\textit{system boundary}) separa el sistema bajo desarrollo de su contexto externo, existiendo una \textbf{zona gris} que se refina conforme avanza la elicitación \cite{pohl2025}. Para analizar sistemáticamente dicho entorno sin omitir restricciones, Pohl establece \textbf{ocho perspectivas de contexto}: (1) Perspectiva de Sujeto/Datos, (2) Perspectiva de Uso, (3) Perspectiva de Entorno TI, (4) Perspectiva de Desarrollo, (5) Perspectiva Física, (6) Perspectiva Social/Cultural, (7) Perspectiva Legal/Regulatoria y (8) Perspectiva Económica/Organizacional \cite{pohl2025}.
\textbf{Explicación Propia:} Muchos proyectos fracasan no por defectos en sus cálculos internos, sino por ignorar las condiciones del entorno donde el sistema será operado. Trazar con precisión el límite del sistema permite acordar con los interesados qué procesos quedan automatizados dentro de la Plataforma de Ventas y qué entidades externas (servicios gubernamentales, pasarelas bancarias, periféricos físicos) solo intercambian información a través de interfaces. Asimismo, recorrer las ocho perspectivas de Pohl actúa como un instrumento metodológico exhaustivo que revela requisitos ocultos, desde la ergonomía física del mostrador hasta las leyes tributarias y la jerga cultural de los compradores.
\textbf{Ejemplo Vinculado al PFC:} Al delimitar la zona gris del PFC, determinamos que la declaración patrimonial y contabilidad general quedan fuera del límite del sistema (solo se emiten y autorizan los comprobantes electrónicos hacia el SRI), mientras que la recomendación técnica de sustitutos y el control de stock sí forman parte del núcleo interno. Asimismo, al aplicar la \textit{Perspectiva Física} de Pohl, identificamos que los mostradores de almacenes eléctricos tienen espacio reducido y polvo, originando el requisito de operar toda la facturación mediante atajos de teclado (F1--F12) y lector de código de barras sin depender del ratón.
\subsection{2.6 Proceso y Framework de Ingeniería de Requisitos (Desarrollado por Jeanpierre Robinson)}
\subsubsection{Concepto 1: Actividades Core del Proceso de Requisitos (Elicitación, Análisis, Especificación y Validación)}
\textbf{Definición Formal:} Tanto en el framework de Pohl (2025) \cite{pohl2025} como en el SWEBOK v4.0 \cite{swebok2024} e ISO/IEC/IEEE 29148:2018 \cite{ieee29148}, el proceso de ingeniería de requisitos se estructura en cuatro actividades fundamentales (\textit{core activities}) interdependientes e iterativas:
\begin{enumerate}
	\item \textbf{Elicitación (Descubrimiento):} Identificación de fuentes y extracción de necesidades de los stakeholders mediante entrevistas, observación en sitio, talleres o análisis documental \cite{pohl2025, wiegers2013}.
	\item \textbf{Análisis y Negociación:} Exploración de factibilidad, detección de conflictos entre interesados, evaluación de riesgos y consenso de prioridades \cite{swebok2024, ireb2023}.
	\item \textbf{Especificación y Documentación:} Registro estructurado de los requisitos acordados en artefactos formales (SRS, modelos conceptuales, tablas EARS) siguiendo normas estandarizadas \cite{ieee29148, mavin2009ears}.
	\item \textbf{Validación y Verificación:} Comprobación de que los requisitos documentados reflejan exactamente lo que el cliente necesita (validación) y cumplen con las propiedades de calidad normativa (verificación) \cite{ieee29148, pohl2025}.
\end{enumerate}
\textbf{Explicación Propia:} Estas cuatro actividades no conforman una secuencia lineal rígida, sino un ciclo iterativo de refinamiento continuo. Durante la elicitación es habitual recolectar expectativas subjetivas o contradictorias entre el propietario y los operarios; la fase de análisis depura y negocia esas tensiones antes de formalizarlas en la especificación. Finalmente, la validación cierra el ciclo asegurando mediante revisiones formales, prototipos y protocolos de aceptación que no exista distorsión entre el problema real del negocio y la especificación técnica acordada.
\textbf{Ejemplo Vinculado al PFC:} En el PFC ejecutamos el ciclo completo: (1) en \textit{Elicitación}, realizamos observación directa en mostrador y entrevistas al cajero y administrador en Quevedo; (2) en \textit{Análisis}, resolvimos el conflicto entre presupuesto y velocidad acordando el uso de un servicio externo de inferencia IA acotado al catálogo local; (3) en \textit{Especificación}, documentamos los 15 requisitos bajo sintaxis EARS e ISO/IEC 25010; y (4) en \textit{Validación}, verificamos cada requisito contra las propiedades de IEEE 29148 y definimos pruebas de aceptación con escenarios reales de productos eléctricos agotados.
\subsubsection{Concepto 2: Actividades Transversales (Gestión de Requisitos y Trazabilidad) y Actores del Proceso}
\textbf{Definición Formal:} De acuerdo con IREB CPRE v3.1 \cite{ireb2023}, Pohl \cite{pohl2025} y PMBOK 7.ª ed. \cite{pmbok}, la \textbf{Gestión de Requisitos} es una actividad transversal continua que abarca el control de versiones, la administración de líneas base (\textit{baselines}), el control de cambios y el mantenimiento de la matriz de trazabilidad a lo largo del ciclo de vida. En este proceso participan actores con roles específicos: el \textbf{Ingeniero de Requisitos / Analista}, el \textbf{Cliente / Patrocinador}, el \textbf{Usuario Operativo}, el \textbf{Arquitecto / Desarrollador}, el \textbf{Especialista de Calidad (QA)} y los \textbf{Entes Reguladores} \cite{ieee29148, wiegers2013}.
\textbf{Explicación Propia:} Los requisitos evolucionan conforme cambian las condiciones del mercado o las normativas legales; sin una gestión transversal disciplinada, cualquier proyecto sufre desviaciones de alcance (\textit{scope creep}) y pérdida de consistencia documental. Del mismo modo, delimitar con precisión las responsabilidades de cada actor evita vacíos de autoridad: el patrocinador prioriza el valor de negocio, el usuario valida la operabilidad diaria, el ingeniero de requisitos custodia la calidad y trazabilidad de la SRS, y el equipo técnico evalúa la factibilidad constructiva.
\textbf{Ejemplo Vinculado al PFC:} En nuestro proyecto establecimos una línea base formal de requisitos (\textit{Baseline v1.0}) bajo control de versiones documental. Cualquier solicitud de cambio propuesta por el Propietario (por ejemplo, incorporar pagos mediante billeteras digitales QR) es evaluada por el Analista Líder (José Gilces) utilizando la matriz de trazabilidad bidireccional para identificar qué requisitos funcionales, de interfaz y de seguridad se ven impactados antes de su aprobación.
